\documentclass[11pt,a4paper]{article}
\newcommand{\gradelabel}{Grade 11}
\newcommand{\leconhdr}{Grandeurs physiques et signaux}
% Préambule commun aux plans de cours de robotique (Grades 7 à 12).
% Chaque plan de cours commence par \documentclass[11pt,a4paper]{article}, définit
% \gradelabel et \leconhdr (en-têtes de page), puis % Préambule commun aux plans de cours de robotique (Grades 7 à 12).
% Chaque plan de cours commence par \documentclass[11pt,a4paper]{article}, définit
% \gradelabel et \leconhdr (en-têtes de page), puis % Préambule commun aux plans de cours de robotique (Grades 7 à 12).
% Chaque plan de cours commence par \documentclass[11pt,a4paper]{article}, définit
% \gradelabel et \leconhdr (en-têtes de page), puis \input{../../preambule_robotique.tex}.
\usepackage[T1]{fontenc}
\usepackage[utf8]{inputenc}
\usepackage{lmodern}
\usepackage[french]{babel}
\usepackage[margin=2cm,headheight=14pt]{geometry}
\usepackage[most]{tcolorbox}
\usepackage{xcolor,tabularx,booktabs,array,enumitem,fancyhdr,amssymb,tikz}
\usetikzlibrary{arrows.meta,positioning}

\pagestyle{fancy}
\fancyhf{}
\lhead{\small Robotique -- \gradelabel{} -- Novembre 2026}
\rhead{\small \leconhdr}
\cfoot{\thepage}
\setlist{itemsep=1pt,topsep=2pt,leftmargin=*}
\renewcommand{\arraystretch}{1.35}
\newcolumntype{Y}{>{\raggedright\arraybackslash}X}
\newcommand{\cb}{$\square$\,}
\newcommand{\fl}{$\rightarrow$}
\newcommand{\lig}{\noindent\rule{\linewidth}{0.3pt}\\[5mm]}

\newtcolorbox{etape}[3]{enhanced,breakable,colback=#1!6,colframe=#1!70!black,
  coltitle=white,fonttitle=\bfseries,title={#2\hfill #3},boxrule=0.8pt,arc=2mm,
  before skip=8pt,after skip=8pt}
\newtcolorbox{astuce}[1][Différenciation et besoins particuliers]{breakable,
  colback=yellow!12,colframe=yellow!50!black,fonttitle=\bfseries\small,
  title={#1},boxrule=0.5pt,arc=1mm,before skip=6pt,after skip=6pt}
\newtcolorbox{code}[1][Code]{breakable,colback=black!4,colframe=black!60,
  fonttitle=\bfseries\small,title={#1},boxrule=0.5pt,arc=1mm,fontupper=\ttfamily\small,
  before skip=6pt,after skip=6pt}
% Blocs de code coloré (écrire le code sans accents : comme dans l'éditeur)
\newtcblisting{arduino}[1][Code Arduino]{listing only,breakable,colback=black!4,
  colframe=black!60,fonttitle=\bfseries\small,title={#1},boxrule=0.5pt,arc=1mm,
  before skip=6pt,after skip=6pt,
  listing options={language=C++,basicstyle=\ttfamily\small,columns=fullflexible,
  keepspaces=true,showstringspaces=false,commentstyle=\color{green!40!black}\itshape,
  keywordstyle=\color{blue!70!black}\bfseries}}
\newtcblisting{python}[1][Code Python]{listing only,breakable,colback=black!4,
  colframe=black!60,fonttitle=\bfseries\small,title={#1},boxrule=0.5pt,arc=1mm,
  before skip=6pt,after skip=6pt,
  listing options={language=Python,basicstyle=\ttfamily\small,columns=fullflexible,
  keepspaces=true,showstringspaces=false,commentstyle=\color{green!40!black}\itshape,
  keywordstyle=\color{blue!70!black}\bfseries}}
\newcommand{\entete}[2]{%
  \begin{tcolorbox}[colback=#1!10,colframe=#1!70!black,boxrule=0.8pt,arc=2mm]
  {\Large\bfseries #2}\hfill Nom : \rule{4cm}{0.3pt}\quad Date : \rule{2.2cm}{0.3pt}
  \end{tcolorbox}}

% Bandeau de titre : #1 titre, #2 semaine
\newcommand{\titreblock}[2]{%
\begin{tcolorbox}[colback=teal!10,colframe=teal!70!black,arc=2mm,boxrule=1pt]
\centering
{\LARGE\bfseries #1}\\[2mm]
{\large Robotique -- \gradelabel{} -- Plan de cours de 2 heures}\\[1mm]
Enseignant : Francis Ki \quad|\quad #2
\end{tcolorbox}}

% Tableau du minutage : 6 contenus (accroche, mini-leçon, guidée, autonome, mise en commun, exit ticket)
\newcommand{\minutage}[6]{%
\begin{tabularx}{\linewidth}{>{\bfseries}p{2.2cm}p{2.6cm}Y p{3cm}}
\toprule
Minutes & Étape & Contenu & Organisation\\
\midrule
0 -- 10 & Accroche & #1 & Classe entière\\
10 -- 30 & Mini-leçon & #2 & Classe entière\\
30 -- 60 & Pratique guidée & #3 & Collectif puis binômes\\
60 -- 95 & Pratique autonome & #4 & Individuel sur poste\\
95 -- 110 & Mise en commun & #5 & Classe entière\\
110 -- 120 & Exit ticket & #6 & Individuel sur poste\\
\bottomrule
\end{tabularx}}

% Étayage progressif : #1 aide 1, #2 aide 2, #3 aide 3
\newcommand{\aides}[3]{%
\textbf{Étayage progressif de l'aide.}
\begin{center}
\begin{tabularx}{\linewidth}{>{\bfseries}p{2.4cm}Y p{4.4cm}}
\toprule
Niveau d'aide & Contenu & Qui intervient\\
\midrule
Aide 1 & #1 & Le binôme seul\\
Aide 2 & #2 & Le binôme avec la carte d'aide\\
Aide 3 & #3 & L'enseignant\\
\bottomrule
\end{tabularx}
\end{center}}

% Les trois groupes de la pratique autonome : #1 tâche A, #2 tâche B, #3 tâche C
\newcommand{\groupes}[3]{%
\textbf{Consigne commune\,:} chaque élève ouvre le fichier de son groupe, le renomme avec son prénom et travaille seul sur son poste. L'enseignant passe d'abord auprès du groupe A pendant 10 min.

\medskip
\begin{tabularx}{\linewidth}{>{\bfseries}p{3cm}Y p{2.2cm}}
\toprule
Groupe & Tâche sur poste & Annexe\\
\midrule
A : en dessous du niveau & #1 & Annexe C\\
B : au niveau & #2 & Annexe D\\
C : au-dessus du niveau & #3 & Annexe E\\
\bottomrule
\end{tabularx}}

% Mise en commun : #1 groupe A, #2 groupe B, #3 groupe C, #4 conclusion
\newcommand{\miseencommun}[4]{%
\begin{etape}{teal}{5. Mise en commun}{95 -- 110 min}
\begin{enumerate}
 \item (2 min) Retour au calme, écrans mis en veille.
 \item (9 min) Trois passages courts de 3 min, avec valorisation :
 \begin{itemize}
  \item \textbf{Groupe A} : #1
  \item \textbf{Groupe B} : #2
  \item \textbf{Groupe C} : #3
 \end{itemize}
 \item (4 min) Conclusion : #4
\end{enumerate}
\end{etape}}

% Clôture : #1 question 1, #2 question 2
\newcommand{\cloture}[2]{%
\begin{etape}{gray}{6. Clôture : exit ticket}{110 -- 120 min}
Formulaire en ligne ou papier (annexe G) :
\begin{enumerate}
 \item #1
 \item #2
 \item Auto-évaluation : \emph{Je maîtrise / Je progresse / J'ai besoin d'aide.}
\end{enumerate}
Lire les réponses avant la séance suivante pour ajuster les groupes de besoin.
\end{etape}}

% Grille d'observation : #1 #2 #3 intitulés de colonnes
\newcommand{\grille}[3]{%
\subsection*{Grille d'observation (enseignant)}
\begin{tabularx}{\linewidth}{Y p{2.4cm} p{2.4cm} p{2.4cm} Y}
\toprule
Élève & #1 & #2 & #3 & Groupe suivant\\
\midrule
 & & & & \\ & & & & \\ & & & & \\ & & & & \\ & & & & \\
\bottomrule
\end{tabularx}}

% Signature de l'exit ticket
\newcommand{\autoeval}{Je me sens : \cb Je maîtrise \quad \cb Je progresse \quad \cb J'ai besoin d'aide}
.
\usepackage[T1]{fontenc}
\usepackage[utf8]{inputenc}
\usepackage{lmodern}
\usepackage[french]{babel}
\usepackage[margin=2cm,headheight=14pt]{geometry}
\usepackage[most]{tcolorbox}
\usepackage{xcolor,tabularx,booktabs,array,enumitem,fancyhdr,amssymb,tikz}
\usetikzlibrary{arrows.meta,positioning}

\pagestyle{fancy}
\fancyhf{}
\lhead{\small Robotique -- \gradelabel{} -- Novembre 2026}
\rhead{\small \leconhdr}
\cfoot{\thepage}
\setlist{itemsep=1pt,topsep=2pt,leftmargin=*}
\renewcommand{\arraystretch}{1.35}
\newcolumntype{Y}{>{\raggedright\arraybackslash}X}
\newcommand{\cb}{$\square$\,}
\newcommand{\fl}{$\rightarrow$}
\newcommand{\lig}{\noindent\rule{\linewidth}{0.3pt}\\[5mm]}

\newtcolorbox{etape}[3]{enhanced,breakable,colback=#1!6,colframe=#1!70!black,
  coltitle=white,fonttitle=\bfseries,title={#2\hfill #3},boxrule=0.8pt,arc=2mm,
  before skip=8pt,after skip=8pt}
\newtcolorbox{astuce}[1][Différenciation et besoins particuliers]{breakable,
  colback=yellow!12,colframe=yellow!50!black,fonttitle=\bfseries\small,
  title={#1},boxrule=0.5pt,arc=1mm,before skip=6pt,after skip=6pt}
\newtcolorbox{code}[1][Code]{breakable,colback=black!4,colframe=black!60,
  fonttitle=\bfseries\small,title={#1},boxrule=0.5pt,arc=1mm,fontupper=\ttfamily\small,
  before skip=6pt,after skip=6pt}
% Blocs de code coloré (écrire le code sans accents : comme dans l'éditeur)
\newtcblisting{arduino}[1][Code Arduino]{listing only,breakable,colback=black!4,
  colframe=black!60,fonttitle=\bfseries\small,title={#1},boxrule=0.5pt,arc=1mm,
  before skip=6pt,after skip=6pt,
  listing options={language=C++,basicstyle=\ttfamily\small,columns=fullflexible,
  keepspaces=true,showstringspaces=false,commentstyle=\color{green!40!black}\itshape,
  keywordstyle=\color{blue!70!black}\bfseries}}
\newtcblisting{python}[1][Code Python]{listing only,breakable,colback=black!4,
  colframe=black!60,fonttitle=\bfseries\small,title={#1},boxrule=0.5pt,arc=1mm,
  before skip=6pt,after skip=6pt,
  listing options={language=Python,basicstyle=\ttfamily\small,columns=fullflexible,
  keepspaces=true,showstringspaces=false,commentstyle=\color{green!40!black}\itshape,
  keywordstyle=\color{blue!70!black}\bfseries}}
\newcommand{\entete}[2]{%
  \begin{tcolorbox}[colback=#1!10,colframe=#1!70!black,boxrule=0.8pt,arc=2mm]
  {\Large\bfseries #2}\hfill Nom : \rule{4cm}{0.3pt}\quad Date : \rule{2.2cm}{0.3pt}
  \end{tcolorbox}}

% Bandeau de titre : #1 titre, #2 semaine
\newcommand{\titreblock}[2]{%
\begin{tcolorbox}[colback=teal!10,colframe=teal!70!black,arc=2mm,boxrule=1pt]
\centering
{\LARGE\bfseries #1}\\[2mm]
{\large Robotique -- \gradelabel{} -- Plan de cours de 2 heures}\\[1mm]
Enseignant : Francis Ki \quad|\quad #2
\end{tcolorbox}}

% Tableau du minutage : 6 contenus (accroche, mini-leçon, guidée, autonome, mise en commun, exit ticket)
\newcommand{\minutage}[6]{%
\begin{tabularx}{\linewidth}{>{\bfseries}p{2.2cm}p{2.6cm}Y p{3cm}}
\toprule
Minutes & Étape & Contenu & Organisation\\
\midrule
0 -- 10 & Accroche & #1 & Classe entière\\
10 -- 30 & Mini-leçon & #2 & Classe entière\\
30 -- 60 & Pratique guidée & #3 & Collectif puis binômes\\
60 -- 95 & Pratique autonome & #4 & Individuel sur poste\\
95 -- 110 & Mise en commun & #5 & Classe entière\\
110 -- 120 & Exit ticket & #6 & Individuel sur poste\\
\bottomrule
\end{tabularx}}

% Étayage progressif : #1 aide 1, #2 aide 2, #3 aide 3
\newcommand{\aides}[3]{%
\textbf{Étayage progressif de l'aide.}
\begin{center}
\begin{tabularx}{\linewidth}{>{\bfseries}p{2.4cm}Y p{4.4cm}}
\toprule
Niveau d'aide & Contenu & Qui intervient\\
\midrule
Aide 1 & #1 & Le binôme seul\\
Aide 2 & #2 & Le binôme avec la carte d'aide\\
Aide 3 & #3 & L'enseignant\\
\bottomrule
\end{tabularx}
\end{center}}

% Les trois groupes de la pratique autonome : #1 tâche A, #2 tâche B, #3 tâche C
\newcommand{\groupes}[3]{%
\textbf{Consigne commune\,:} chaque élève ouvre le fichier de son groupe, le renomme avec son prénom et travaille seul sur son poste. L'enseignant passe d'abord auprès du groupe A pendant 10 min.

\medskip
\begin{tabularx}{\linewidth}{>{\bfseries}p{3cm}Y p{2.2cm}}
\toprule
Groupe & Tâche sur poste & Annexe\\
\midrule
A : en dessous du niveau & #1 & Annexe C\\
B : au niveau & #2 & Annexe D\\
C : au-dessus du niveau & #3 & Annexe E\\
\bottomrule
\end{tabularx}}

% Mise en commun : #1 groupe A, #2 groupe B, #3 groupe C, #4 conclusion
\newcommand{\miseencommun}[4]{%
\begin{etape}{teal}{5. Mise en commun}{95 -- 110 min}
\begin{enumerate}
 \item (2 min) Retour au calme, écrans mis en veille.
 \item (9 min) Trois passages courts de 3 min, avec valorisation :
 \begin{itemize}
  \item \textbf{Groupe A} : #1
  \item \textbf{Groupe B} : #2
  \item \textbf{Groupe C} : #3
 \end{itemize}
 \item (4 min) Conclusion : #4
\end{enumerate}
\end{etape}}

% Clôture : #1 question 1, #2 question 2
\newcommand{\cloture}[2]{%
\begin{etape}{gray}{6. Clôture : exit ticket}{110 -- 120 min}
Formulaire en ligne ou papier (annexe G) :
\begin{enumerate}
 \item #1
 \item #2
 \item Auto-évaluation : \emph{Je maîtrise / Je progresse / J'ai besoin d'aide.}
\end{enumerate}
Lire les réponses avant la séance suivante pour ajuster les groupes de besoin.
\end{etape}}

% Grille d'observation : #1 #2 #3 intitulés de colonnes
\newcommand{\grille}[3]{%
\subsection*{Grille d'observation (enseignant)}
\begin{tabularx}{\linewidth}{Y p{2.4cm} p{2.4cm} p{2.4cm} Y}
\toprule
Élève & #1 & #2 & #3 & Groupe suivant\\
\midrule
 & & & & \\ & & & & \\ & & & & \\ & & & & \\ & & & & \\
\bottomrule
\end{tabularx}}

% Signature de l'exit ticket
\newcommand{\autoeval}{Je me sens : \cb Je maîtrise \quad \cb Je progresse \quad \cb J'ai besoin d'aide}
.
\usepackage[T1]{fontenc}
\usepackage[utf8]{inputenc}
\usepackage{lmodern}
\usepackage[french]{babel}
\usepackage[margin=2cm,headheight=14pt]{geometry}
\usepackage[most]{tcolorbox}
\usepackage{xcolor,tabularx,booktabs,array,enumitem,fancyhdr,amssymb,tikz}
\usetikzlibrary{arrows.meta,positioning}

\pagestyle{fancy}
\fancyhf{}
\lhead{\small Robotique -- \gradelabel{} -- Novembre 2026}
\rhead{\small \leconhdr}
\cfoot{\thepage}
\setlist{itemsep=1pt,topsep=2pt,leftmargin=*}
\renewcommand{\arraystretch}{1.35}
\newcolumntype{Y}{>{\raggedright\arraybackslash}X}
\newcommand{\cb}{$\square$\,}
\newcommand{\fl}{$\rightarrow$}
\newcommand{\lig}{\noindent\rule{\linewidth}{0.3pt}\\[5mm]}

\newtcolorbox{etape}[3]{enhanced,breakable,colback=#1!6,colframe=#1!70!black,
  coltitle=white,fonttitle=\bfseries,title={#2\hfill #3},boxrule=0.8pt,arc=2mm,
  before skip=8pt,after skip=8pt}
\newtcolorbox{astuce}[1][Différenciation et besoins particuliers]{breakable,
  colback=yellow!12,colframe=yellow!50!black,fonttitle=\bfseries\small,
  title={#1},boxrule=0.5pt,arc=1mm,before skip=6pt,after skip=6pt}
\newtcolorbox{code}[1][Code]{breakable,colback=black!4,colframe=black!60,
  fonttitle=\bfseries\small,title={#1},boxrule=0.5pt,arc=1mm,fontupper=\ttfamily\small,
  before skip=6pt,after skip=6pt}
% Blocs de code coloré (écrire le code sans accents : comme dans l'éditeur)
\newtcblisting{arduino}[1][Code Arduino]{listing only,breakable,colback=black!4,
  colframe=black!60,fonttitle=\bfseries\small,title={#1},boxrule=0.5pt,arc=1mm,
  before skip=6pt,after skip=6pt,
  listing options={language=C++,basicstyle=\ttfamily\small,columns=fullflexible,
  keepspaces=true,showstringspaces=false,commentstyle=\color{green!40!black}\itshape,
  keywordstyle=\color{blue!70!black}\bfseries}}
\newtcblisting{python}[1][Code Python]{listing only,breakable,colback=black!4,
  colframe=black!60,fonttitle=\bfseries\small,title={#1},boxrule=0.5pt,arc=1mm,
  before skip=6pt,after skip=6pt,
  listing options={language=Python,basicstyle=\ttfamily\small,columns=fullflexible,
  keepspaces=true,showstringspaces=false,commentstyle=\color{green!40!black}\itshape,
  keywordstyle=\color{blue!70!black}\bfseries}}
\newcommand{\entete}[2]{%
  \begin{tcolorbox}[colback=#1!10,colframe=#1!70!black,boxrule=0.8pt,arc=2mm]
  {\Large\bfseries #2}\hfill Nom : \rule{4cm}{0.3pt}\quad Date : \rule{2.2cm}{0.3pt}
  \end{tcolorbox}}

% Bandeau de titre : #1 titre, #2 semaine
\newcommand{\titreblock}[2]{%
\begin{tcolorbox}[colback=teal!10,colframe=teal!70!black,arc=2mm,boxrule=1pt]
\centering
{\LARGE\bfseries #1}\\[2mm]
{\large Robotique -- \gradelabel{} -- Plan de cours de 2 heures}\\[1mm]
Enseignant : Francis Ki \quad|\quad #2
\end{tcolorbox}}

% Tableau du minutage : 6 contenus (accroche, mini-leçon, guidée, autonome, mise en commun, exit ticket)
\newcommand{\minutage}[6]{%
\begin{tabularx}{\linewidth}{>{\bfseries}p{2.2cm}p{2.6cm}Y p{3cm}}
\toprule
Minutes & Étape & Contenu & Organisation\\
\midrule
0 -- 10 & Accroche & #1 & Classe entière\\
10 -- 30 & Mini-leçon & #2 & Classe entière\\
30 -- 60 & Pratique guidée & #3 & Collectif puis binômes\\
60 -- 95 & Pratique autonome & #4 & Individuel sur poste\\
95 -- 110 & Mise en commun & #5 & Classe entière\\
110 -- 120 & Exit ticket & #6 & Individuel sur poste\\
\bottomrule
\end{tabularx}}

% Étayage progressif : #1 aide 1, #2 aide 2, #3 aide 3
\newcommand{\aides}[3]{%
\textbf{Étayage progressif de l'aide.}
\begin{center}
\begin{tabularx}{\linewidth}{>{\bfseries}p{2.4cm}Y p{4.4cm}}
\toprule
Niveau d'aide & Contenu & Qui intervient\\
\midrule
Aide 1 & #1 & Le binôme seul\\
Aide 2 & #2 & Le binôme avec la carte d'aide\\
Aide 3 & #3 & L'enseignant\\
\bottomrule
\end{tabularx}
\end{center}}

% Les trois groupes de la pratique autonome : #1 tâche A, #2 tâche B, #3 tâche C
\newcommand{\groupes}[3]{%
\textbf{Consigne commune\,:} chaque élève ouvre le fichier de son groupe, le renomme avec son prénom et travaille seul sur son poste. L'enseignant passe d'abord auprès du groupe A pendant 10 min.

\medskip
\begin{tabularx}{\linewidth}{>{\bfseries}p{3cm}Y p{2.2cm}}
\toprule
Groupe & Tâche sur poste & Annexe\\
\midrule
A : en dessous du niveau & #1 & Annexe C\\
B : au niveau & #2 & Annexe D\\
C : au-dessus du niveau & #3 & Annexe E\\
\bottomrule
\end{tabularx}}

% Mise en commun : #1 groupe A, #2 groupe B, #3 groupe C, #4 conclusion
\newcommand{\miseencommun}[4]{%
\begin{etape}{teal}{5. Mise en commun}{95 -- 110 min}
\begin{enumerate}
 \item (2 min) Retour au calme, écrans mis en veille.
 \item (9 min) Trois passages courts de 3 min, avec valorisation :
 \begin{itemize}
  \item \textbf{Groupe A} : #1
  \item \textbf{Groupe B} : #2
  \item \textbf{Groupe C} : #3
 \end{itemize}
 \item (4 min) Conclusion : #4
\end{enumerate}
\end{etape}}

% Clôture : #1 question 1, #2 question 2
\newcommand{\cloture}[2]{%
\begin{etape}{gray}{6. Clôture : exit ticket}{110 -- 120 min}
Formulaire en ligne ou papier (annexe G) :
\begin{enumerate}
 \item #1
 \item #2
 \item Auto-évaluation : \emph{Je maîtrise / Je progresse / J'ai besoin d'aide.}
\end{enumerate}
Lire les réponses avant la séance suivante pour ajuster les groupes de besoin.
\end{etape}}

% Grille d'observation : #1 #2 #3 intitulés de colonnes
\newcommand{\grille}[3]{%
\subsection*{Grille d'observation (enseignant)}
\begin{tabularx}{\linewidth}{Y p{2.4cm} p{2.4cm} p{2.4cm} Y}
\toprule
Élève & #1 & #2 & #3 & Groupe suivant\\
\midrule
 & & & & \\ & & & & \\ & & & & \\ & & & & \\ & & & & \\
\bottomrule
\end{tabularx}}

% Signature de l'exit ticket
\newcommand{\autoeval}{Je me sens : \cb Je maîtrise \quad \cb Je progresse \quad \cb J'ai besoin d'aide}


\begin{document}
\titreblock{Grandeurs physiques et signaux}{Semaine du 23 au 27 novembre 2026}

%==================================================================
\section{Informations générales}
\begin{tabularx}{\linewidth}{>{\bfseries}p{3.6cm}Y}
\toprule
Titre & Grandeurs physiques et signaux\\
Classe / effectif & Grade 11, 11 à 20 élèves, niveaux hétérogènes, élèves à besoins particuliers\\
Durée & 2 heures (120 min)\\
Lieu & Salle informatique, un poste par élève\\
Matériel et logiciels &
\begin{itemize}
 \item Un poste par élève avec navigateur web.
 \item \textbf{Wokwi} (wokwi.com, gratuit) : projet ESP32 modèle avec une thermistance (capteur de température NTC) sur GPIO 34 et une photorésistance (capteur de lumière) sur GPIO 35. Le code est fourni.
 \item Un tableur (Google Sheets ou LibreOffice Calc) pour les graphiques.
 \item Vidéoprojecteur ou tableau interactif.
\end{itemize}\\
Pré-requis &
Séances précédentes : chaîne d'information, capteurs analogiques et numériques, CAN 12 bits (0 à 4095 pour 0 à 3,3~V).\\
\bottomrule
\end{tabularx}

\begin{astuce}[Préparation avant le cours (environ 45 min)]
\begin{enumerate}
 \item Créer le projet Wokwi avec les deux capteurs et coller le code de l'annexe A. Vérifier les noms et les broches des composants dans votre version du simulateur.
 \item Préparer un tableur modèle avec les colonnes (temps, valeur, tension).
 \item Préparer un graphique d'exemple de température sur 24~h pour l'accroche.
 \item Préparer les fichiers des annexes C à E et le formulaire de sortie.
\end{enumerate}
\end{astuce}

%==================================================================
\section{Objectifs d'apprentissage}
\begin{tcolorbox}[colback=gray!8,colframe=gray!60!black,arc=1mm,boxrule=0.5pt]
\textbf{Objectifs du programme}
\begin{itemize}
 \item Identifier une grandeur mesurée.
 \item Interpréter un signal simple.
\end{itemize}
\textbf{À la fin de la séance, l'élève est capable\,:}
\begin{itemize}
 \item d'associer à un capteur la grandeur physique qu'il mesure et son unité\,;
 \item de lire un signal (courbe d'une valeur en fonction du temps) et d'en dégager un minimum, un maximum et une tendance\,;
 \item de choisir un seuil à partir d'un signal pour déclencher une action.
\end{itemize}
\textbf{Compétences numériques\,:} relever des mesures, tracer et lire une courbe, justifier une décision par les données.
\end{tcolorbox}

%==================================================================
\section{Déroulé minuté (120 min)}
\minutage{Que raconte cette courbe de température ?}{Grandeur, unité, signal, seuil}{Lecture d'un signal en collectif, puis relevé et courbe en binômes}{Trois niveaux de besoin}{Valoriser chaque groupe}{Trois questions sur poste}

%------------------------------------------------------------------
\begin{etape}{orange}{1. Accroche : que raconte cette courbe ?}{0 -- 10 min}
\begin{enumerate}
 \item (5 min) Projeter la courbe de la température d'une journée (24 h) : montée le matin, maximum l'après-midi, descente le soir. \emph{À quelle heure fait-il le plus chaud ? Quand le chauffage doit-il se déclencher ?}
 \item (5 min) Question-problème : \textbf{«~Comment un système peut-il comprendre ce que raconte une courbe et décider tout seul ?~»}
\end{enumerate}
\end{etape}

%------------------------------------------------------------------
\begin{etape}{blue}{2. Mini-leçon : grandeurs, signaux, seuils}{10 -- 30 min}
\textbf{Support\,:} annexe A. Quatre temps de 5 min.

\textbf{Temps 1 : la grandeur physique.} C'est ce qu'on mesure. Elle a une unité.
\begin{center}
\begin{tabularx}{\linewidth}{>{\bfseries}p{3.4cm}Y Y}
\toprule
Grandeur & Unité & Capteur possible\\
\midrule
Température & degré Celsius (°C) & Thermistance, sonde\\
Luminosité & lux (lx) & Photorésistance\\
Distance & mètre (m) & Capteur ultrason\\
Vitesse & mètre par seconde (m/s) & Codeur\\
Humidité & pour cent (\%) & Capteur d'humidité\\
\bottomrule
\end{tabularx}
\end{center}

\textbf{Temps 2 : un signal est une grandeur qui varie.} On le représente par une courbe : le temps en abscisse, la valeur en ordonnée.

\textbf{Temps 3 : lire un signal.} Repérer le minimum, le maximum, la tendance (monte, descend, stable), les valeurs anormales (pics). Le bruit est une petite variation parasite.

\begin{center}
\begin{tikzpicture}[scale=0.9]
\draw[->] (0,0) -- (8,0) node[right]{\small temps};
\draw[->] (0,0) -- (0,3.2) node[above]{\small valeur};
\draw[thick,blue!70!black] (0,0.8) .. controls (2,0.6) and (3,3) .. (4,2.8) .. controls (5,2.6) and (6,0.7) .. (8,0.9);
\draw[dashed,red!70!black] (0,1.6) -- (8,1.6) node[right,font=\small]{seuil};
\end{tikzpicture}
\end{center}

\textbf{Temps 4 : le seuil.} Un seuil est une valeur qui déclenche une action. On le choisit à partir du signal, en tenant compte du bruit : un seuil trop proche d'une valeur courante fait «~clignoter~» la commande.

\begin{astuce}
Utiliser des couleurs pour le minimum, le maximum et le seuil sur toutes les courbes. Fournir des courbes déjà tracées pour les élèves qui ont du mal à construire un graphique. Lire à voix haute les axes avant d'interpréter.
\end{astuce}
\end{etape}

%------------------------------------------------------------------
\begin{etape}{green}{3. Pratique guidée}{30 -- 60 min}
\textbf{Fiche\,:} annexe B. Les élèves travaillent en binômes devant leur poste.

\textbf{3a. Lecture d'un signal en collectif (12 min).} À partir de la courbe de température de l'accroche, repérer ensemble le minimum, le maximum, l'heure du maximum et la tendance. Placer un seuil de chauffage et indiquer quand il déclenche.

\textbf{3b. Relevé et courbe en binômes (18 min).} Dans Wokwi, faire varier la température de la thermistance (valeurs \(-10\), 0, 10, 20, 30, 40~°C), relever la valeur lue, la reporter dans le tableur et tracer la courbe.

\aides{La courbe de température de l'accroche annotée}{Questions guides : \emph{Que mesure-t-on en ordonnée ? En abscisse ? Quelle est la tendance ? Où mettrais-tu un seuil ?}}{L'enseignant relève un point du tableau avec le binôme, qui poursuit}

\begin{astuce}
Donner le tableau avec les températures déjà écrites et seulement les valeurs lues à compléter. Permettre de dicter les valeurs à son partenaire qui les saisit.
\end{astuce}
\end{etape}

%------------------------------------------------------------------
\begin{etape}{violet}{4. Pratique autonome différenciée}{60 -- 95 min}
\groupes{Associer cinq capteurs à la grandeur mesurée et à l'unité (banque de mots), puis lire une courbe imprimée de température en répondant à quatre questions guidées (minimum, maximum, heure du maximum, tendance).}{Relever dans Wokwi les valeurs de la thermistance pour six températures, tracer la courbe dans le tableur, interpréter (la relation est-elle linéaire ?) et choisir un seuil pour un ventilateur.}{Étudier la photorésistance : relever sept points (de l'obscurité à la pleine lumière), comparer avec la thermistance, choisir un seuil pour un éclairage automatique en tenant compte du bruit, et proposer deux seuils (hystérésis) pour éviter le clignotement.}

\medskip
\textbf{Plan B si Wokwi est indisponible\,:} utiliser le tableau de mesures imprimé et le tableur.
\end{etape}

%------------------------------------------------------------------
\miseencommun{un élève lit la courbe imprimée et dit le maximum et l'heure. Valoriser la lecture précise.}{un binôme présente sa courbe et son seuil. Valoriser le lien entre les données et la décision.}{un élève présente ses deux seuils et explique le risque de clignotement. Valoriser l'analyse.}{\textit{un capteur transforme une grandeur physique en signal ; on lit ce signal pour choisir un seuil et déclencher une action.}}

\cloture{Quelle grandeur mesure une photorésistance et dans quelle unité l'exprime-t-on ?}{Pourquoi faut-il tenir compte du bruit quand on choisit un seuil ?}

%==================================================================
% ANNEXES
%==================================================================
\newpage
\section*{Annexe A : Support résumé de la mini-leçon}
\entete{blue}{Grandeurs physiques et signaux}

\begin{tcolorbox}[colback=yellow!10,colframe=yellow!50!black]
Une \textbf{grandeur physique} est mesurée avec une \textbf{unité}. Un \textbf{signal} est une grandeur qui varie dans le temps. Un \textbf{seuil} déclenche une action.
\end{tcolorbox}

\begin{tabularx}{\linewidth}{>{\bfseries}p{3.2cm}Y Y}
\toprule
Grandeur & Unité & Capteur\\
\midrule
Température & °C & Thermistance\\
Luminosité & lx & Photorésistance\\
Distance & m & Ultrason\\
Humidité & \% & Capteur d'humidité\\
\bottomrule
\end{tabularx}

\begin{arduino}[Code fourni (Wokwi)]
void setup() {
  Serial.begin(115200);
}

void loop() {
  int temp = analogRead(34);    // thermistance
  int lum  = analogRead(35);    // photoresistance
  Serial.print("temp = ");
  Serial.print(temp);
  Serial.print("   lum = ");
  Serial.println(lum);
  delay(500);
}
\end{arduino}

%==================================================================
\newpage
\section*{Annexe B : Fiche de pratique guidée}
\entete{green}{Pratique guidée : lire un signal}

\subsection*{3a. Courbe de température sur 24 h}
Minimum : \dotfill{} °C à \dotfill{} h\quad Maximum : \dotfill{} °C à \dotfill{} h\\
Tendance entre 6 h et 14 h : \dotfill\\
Mon seuil de chauffage : \dotfill{} °C. Il se déclenche entre \dotfill{} h et \dotfill{} h.

\subsection*{3b. Relevé de la thermistance}
\begin{tabularx}{\linewidth}{>{\bfseries}p{3.4cm}Y Y}
\toprule
Température (°C) & Valeur lue (0 à 4095) & Tension (V)\\
\midrule
$-10$ & & \\ 0 & & \\ 10 & & \\ 20 & & \\ 30 & & \\ 40 & & \\
\bottomrule
\end{tabularx}

\begin{tcolorbox}[colback=yellow!10,colframe=yellow!50!black,title={Cartes d'aide à utiliser dans l'ordre},fonttitle=\bfseries\small]
\textbf{Aide 1 :} regarde la courbe annotée de l'accroche.\\
\textbf{Aide 2 :} Que mesure-t-on en ordonnée ? En abscisse ? Quelle est la tendance ? Où mettrais-tu un seuil ?\\
\textbf{Aide 3 :} demande à l'enseignant de relever un point avec toi.
\end{tcolorbox}

%==================================================================
\newpage
\section*{Annexe C : Niveau A (en dessous du niveau)}
\entete{orange}{Niveau A : grandeur, unité, courbe}

\textbf{Banque :} température, luminosité, distance, humidité, °C, lx, m, \%.

\begin{tabularx}{\linewidth}{>{\bfseries}p{4cm}Y Y}
\toprule
Capteur & Grandeur mesurée & Unité\\
\midrule
Thermistance & & \\
Photorésistance & & \\
Capteur ultrason & & \\
Capteur d'humidité & & \\
\bottomrule
\end{tabularx}

\subsection*{Lis la courbe}
\begin{enumerate}
 \item Quel est le maximum ? \dotfill
 \item À quelle heure ? \dotfill
 \item La courbe monte ou descend le matin ? \dotfill
 \item Entoure le moment où le chauffage devrait se déclencher.
\end{enumerate}

%==================================================================
\newpage
\section*{Annexe D : Niveau B (au niveau)}
\entete{blue}{Niveau B : relevé, courbe et seuil}

\begin{tabularx}{\linewidth}{>{\bfseries}p{3.4cm}Y Y}
\toprule
Température (°C) & Valeur lue & Tension (V)\\
\midrule
$-10$ & & \\ 0 & & \\ 10 & & \\ 20 & & \\ 30 & & \\ 40 & & \\
\bottomrule
\end{tabularx}

Trace la courbe (température en abscisse, valeur en ordonnée).

\subsection*{Questions}
\begin{enumerate}
 \item La courbe est-elle une droite ? \dotfill
 \item Quelle valeur lue correspond à 25 °C d'après ta courbe ? \dotfill
 \item Choisis un seuil pour un ventilateur qui démarre à 28 °C : valeur seuil \dotfill
\end{enumerate}

%==================================================================
\newpage
\section*{Annexe E : Niveau C (au-dessus du niveau)}
\entete{violet}{Niveau C : l'éclairage automatique}

\begin{tabularx}{\linewidth}{>{\bfseries}p{3.4cm}Y Y}
\toprule
Éclairement (position) & Valeur lue (photorésistance) & Commentaire\\
\midrule
Obscurité & & \\ Faible & & \\ Moyen & & \\ Fort & & \\ Très fort & & \\
\bottomrule
\end{tabularx}

\begin{enumerate}
 \item Choisis un seuil et justifie : \lig
 \item Explique ce qu'est le bruit et pourquoi la lampe peut clignoter autour du seuil : \lig \lig
 \item Propose deux seuils (allumer, éteindre) pour l'éviter : \dotfill
\end{enumerate}

\textbf{Défi bonus :} compare la sensibilité de la thermistance et de la photorésistance (quelle variation de valeur pour quelle variation de grandeur).

%==================================================================
\newpage
\section*{Annexe F : Corrections et attendus (enseignant)}

\subsection*{Annexe B}
Courbe type : minimum vers 5 h (environ 14 °C), maximum vers 15 h (environ 30 °C), tendance : monte. Un chauffage se déclenche quand la température passe sous un seuil (par exemple 18 °C), le soir et la nuit. Pour une thermistance NTC, la résistance diminue quand la température augmente ; le sens de variation de la valeur lue dépend du montage du module dans le simulateur (\textbf{à vérifier avant le cours}).

\subsection*{Annexe C}
Température (°C) ; luminosité (lx) ; distance (m) ; humidité (\%).

\subsection*{Annexe D}
La courbe d'une thermistance NTC n'est pas une droite. Interpolation : lire sur la courbe. Seuil : valeur lue pour 28 °C.

\subsection*{Annexe E}
Seuil indicatif : environ 1000 à 1500 sur 4095, selon le simulateur. Bruit : petites variations aléatoires qui font franchir le seuil. Hystérésis : allumer sous 900, éteindre au-dessus de 1100.\\
\textbf{Critères de réussite :} lecture correcte de la courbe, seuil justifié, explication du bruit.

%==================================================================
\newpage
\section*{Annexe G : Exit ticket et grille d'observation}
\entete{gray}{Exit ticket : grandeurs et signaux}

\begin{enumerate}
 \item Quelle grandeur mesure une photorésistance et dans quelle unité l'exprime-t-on ?\\ \lig
 \item Pourquoi faut-il tenir compte du bruit quand on choisit un seuil ?\\ \lig
 \item \autoeval
\end{enumerate}
\textbf{Attendus :} 1. La luminosité, en lux. 2. Le bruit fait varier la mesure autour du seuil et la commande peut se déclencher puis s'arrêter sans cesse.

\bigskip
\grille{Grandeur et unité}{Courbe lue}{Seuil justifié}

\end{document}
